% © 2026 Ing. David Jaroš — CC BY-NC-ND 4.0
%
% File: research_tracks/EW/weinberg_angle_odd_spinor_threshold_beta_theorem.tex
% Status: RESEARCH TRACK / BRIDGE-PROOF.
% Purpose: Derive the twisted odd-spinor one-loop beta coefficients used in the Weinberg RG branch.

% UBT-AI-PROVENANCE-BEGIN
% schema: ubt-ai-provenance/v1
% tier: C_working
% ai_assistance: disclosed
% human_review: risk-based
% editorial_responsibility: Ing. David Jaroš
% policy: ../../AI_PROVENANCE.md
% notice: Working material; exhaustive human review is not claimed.
% UBT-AI-PROVENANCE-END

\documentclass[12pt,a4paper]{article}
\usepackage[a4paper,margin=2.5cm]{geometry}
\usepackage{amsmath,amssymb,mathtools,booktabs}
\usepackage[most]{tcolorbox}
\usepackage{hyperref}
\usepackage{xcolor}
\usepackage{microtype}

\hypersetup{colorlinks=true,linkcolor=blue!60!black,urlcolor=blue!60!black,citecolor=green!50!black}

\newtcolorbox{statusbox}[1][]{
  colback=yellow!8!white,
  colframe=orange!70!black,
  arc=3pt,
  boxrule=1pt,
  title=Status,
  #1
}

\title{Odd-Spinor Threshold Beta Coefficients for the Weinberg Angle}
\author{Ing. David Jaroš}
\date{2026-06-14}

\begin{document}
\maketitle

\begin{abstract}
The UBT Weinberg RG branch requires the effective one-loop coefficients
\[
  b_1=\frac{33}{5},\qquad b_2=1,
\]
which are the standard $N=1$ supersymmetric Standard-Model coefficients.
This note derives these coefficients from the UBT odd-winding spinor Layer2
threshold spectrum: three chiral Standard-Model generations plus the conjugate
pair of electroweak doublets $H_u,H_d$ forced by biquaternionic conjugation of
the electroweak readout.  The proof is bridge-level: it is proven relative to
the current UBT bridges that identify odd winding spinors with Berezin measure,
three light winding sectors with the observed generations, and the electroweak
readout with the conjugate doublet pair.
\end{abstract}

\begin{statusbox}
\textbf{Result.}  Under the current UBT odd-spinor Layer2 bridge assumptions,
\[
  (b_1,b_2,b_3)=\left(\frac{33}{5},1,-3\right).
\]
Thus the beta coefficients used in
\texttt{weinberg\_angle\_rg\_twisted\_scale\_closure.tex} are no longer merely
imported as MSSM-like numbers; they follow from the odd-spinor threshold
spectrum.  The remaining Weinberg inputs are the compact scale $M_{\rm UBT}$
and the $Z$-pole electromagnetic coupling obtained from the alpha/QED running
bridge.
\end{statusbox}

\section{Bridge assumptions}

We use four bridge results already present in the repository:
\begin{enumerate}
  \item odd winding modes carry Berezin/Grassmann measure;
  \item the first three light $\psi$-winding sectors encode the three observed
  fermion generations;
  \item the UBT hypercharge bridge gives the standard one-generation
  hypercharge set;
  \item the electroweak readout is closed under biquaternionic conjugation,
  giving a conjugate pair of doublets $H_u,H_d$ rather than a single isolated
  doublet.
\end{enumerate}
The last point is essential: the paired electroweak doublet is the same
structural reason why the twisted branch has the $N=1$ threshold coefficients
rather than the minimal non-supersymmetric SM coefficients.

\section{One-loop formula}

For a twisted odd-spinor threshold with vector multiplets and chiral readout
multiplets, the one-loop coefficients are
\[
  b_i=-3C_2(G_i)+\sum_{\rm chiral}T_i(R).
\]
For the Abelian factor we use the GUT-normalised index
\[
  T_1(R)=\frac35Y^2\dim(SU(3))\dim(SU(2)).
\]
For the non-Abelian fundamental representations,
\[
  T_2(\mathbf 2)=\frac12,
  \qquad
  T_3(\mathbf 3)=\frac12.
\]

\section{Matter contribution}

For three generations of
\[
  Q,\quad u^c,\quad d^c,\quad L,\quad e^c,
\]
the chiral indices are
\[
  \sum_{3\,\mathrm{gen}} T_1=6,
  \qquad
  \sum_{3\,\mathrm{gen}} T_2=6,
  \qquad
  \sum_{3\,\mathrm{gen}} T_3=6.
\]
Explicitly, per generation,
\[
\begin{array}{c|c|c|c|c|c}
\text{field} & SU(3) & SU(2) & Y & T_1 & (T_2,T_3)\\
\hline
Q & 3 & 2 & 1/6 & 1/10 & (3/2,1)\\
u^c & \bar 3 & 1 & -2/3 & 4/5 & (0,1/2)\\
d^c & \bar 3 & 1 & 1/3 & 1/5 & (0,1/2)\\
L & 1 & 2 & -1/2 & 3/10 & (1/2,0)\\
e^c & 1 & 1 & 1 & 3/5 & (0,0)
\end{array}
\]
so one generation contributes $(2,2,2)$ and three generations contribute
$(6,6,6)$.

\section{Electroweak doublet pair}

The conjugate electroweak doublet pair is
\[
  H_u:(1,2,+1/2),
  \qquad
  H_d:(1,2,-1/2).
\]
Each doublet contributes
\[
  T_1=\frac35\left(\frac12\right)^2\cdot2=\frac{3}{10},
  \qquad
  T_2=\frac12,
  \qquad
  T_3=0.
\]
Thus the pair contributes
\[
  (T_1,T_2,T_3)_{H_u+H_d}=\left(\frac35,1,0\right).
\]
The full chiral threshold contribution is therefore
\[
  \sum T_1=\frac{33}{5},
  \qquad
  \sum T_2=7,
  \qquad
  \sum T_3=6.
\]

\section{Vector contribution and final coefficients}

The vector contribution is
\[
  -3C_2(U(1))=0,
  \qquad
  -3C_2(SU(2))=-6,
  \qquad
  -3C_2(SU(3))=-9.
\]
Hence
\[
  b_1=\frac{33}{5},
  \qquad
  b_2=7-6=1,
  \qquad
  b_3=6-9=-3.
\]
This proves the beta coefficients used by the twisted odd-spinor Weinberg RG
branch.

\section{Status}

The theorem closes the beta-coefficient gap conditional on the UBT bridge
assumptions listed above.  It does not by itself derive the compact scale
$M_{\rm UBT}$ or the electromagnetic coupling at the $Z$ pole; those enter the
full low-energy closure theorem.

\end{document}
